\documentclass{article}
\usepackage[T1]{fontenc}
\usepackage{lmodern}
\usepackage{graphicx}
\usepackage{amsmath,amssymb}
\usepackage[a4paper,margin=2cm]{geometry}
\usepackage[absolute,overlay]{textpos}
\setlength{\TPHorizModule}{1mm}
\setlength{\TPVertModule}{1mm}
\usepackage[indonesian]{babel}
\usepackage{hyperref}
\setlength{\emergencystretch}{2em}
% unit: Halaman 91

\begin{document}

% === Halaman 91 (python/native) ===

91

• Contoh: Misalkan kriptanalis menemukan 



cipherteks C dan plainteks berkorepsonden K 



cipherteks E dan plainteks berkoresponden O. 

• Kriptanalis m dan n dari kekongruenan berikut:


2  10m + b (mod 26)  

(i)


4  14m + b (mod 26)  

(ii) 

• Kurangkan (ii) dengan (i), menghasilkan


2  4m (mod 26) 



(iii) 

  Solusi: m = 7

Substitusi m = 7  ke dalam (i),


2  70 + b (mod 26)  


(iv)


Solusi: b = 10.

\end{document}
